\chapter{自指闭环：身份维持、元观察与可验证控制}
\label{chp:topological_grand_unification_of_self_reference}

上一章把宏观潜空间限定为任务相关压缩、受限几何与控制接口，本章进一步回答控制方向从何而来、系统怎样把自身纳入状态、又怎样避免把一次内部一致性误认成现实成功。我们所说的“自我”不是神秘实体、单一存储节点或某个脑区，而是一组跨时间可查询、可修订并参与控制的身份约束；所谓自指也不是逻辑上的无限回望，而是系统把自身状态、历史承诺、能力边界和行动后果重新作为下一轮计算输入的闭环机制。

我们将元观察、历史重入、宏观控制和环境验收分开建模。内部读出可以形成关于“我正在做什么”的估计，却不能仅凭该估计证明内容真实；控制器可以降低内部目标，却不能由此宣布任务完成；生物网络的相关活动可以约束工程假说，却不能把脑区名称直接翻译成软件模块。HSF--HD在这里提供智能一般状态动力学的模型族，AgentOS的$h(t)$、$E$、$\Theta$、$\Psi$、SPC、TCE、CCM、Boundary、Offloading与外部记忆$M_e$仅说明一种可审计实现。以下各节依次定义对象、更新机制、稳定条件、反例和验证协议，使“自我维持”成为可检验的计算命题，而不是由拓扑或热力学词汇担保的结论。

\section{自指对象：身份状态、自我模型与四类功能}
\label{sec:four_element_ontology_topological_basis}

我们先定义离散时间$t_k=k\Delta t$上的联合状态
\begin{equation}
X_k=(h_k,E_k,\Theta_k,M_{e,k},e_k,\nu_k),
\end{equation}
其中$h_k\in\mathbb{R}^{d_h}$是当前工作状态，$E_k$是带时间戳、来源和访问权限的事件记录集合，$\Theta_k=(G_k,\vartheta_k)$包含局部结构图$G_k$与连续参数$\vartheta_k$，$M_{e,k}$是外部记忆引用集合，$e_k$是环境状态，$\nu_k$是版本与制度状态。连续向量无量纲或按接口归一化；时间以秒或任务步计；图编辑、记录提交和权限变化是离散事件，不能假装成连续向量的一部分。任务$q$给出查询族、验收谓词、风险上限和资源预算。系统只能相对于$q$谈“保持了同一个自我”，不存在脱离观察窗口与任务条件的绝对身份向量。

自我模型记为$S_k=(I_k,B_k,C_k,V_k)$。$I_k$是任务相关身份键，包括主体、角色、所有权与责任链；$B_k$是Boundary允许的观察和行动边界；$C_k$是能力、资源、权限和不确定性估计；$V_k$是价值、承诺与优先级的版本化集合。这里的四元组是计算定义，不是四种物理实体。它对应原章的四类功能：身份结构提供跨时索引，元观察器给出关于当前状态的读出，宏观隐状态压缩长期策略方向，执行层把策略翻译为受约束行动。四者可以由同一网络联合实现，也可以分布在多个服务中；理论只要求接口关系和可观察后果，不规定唯一载体。

身份连续性由保持查询而非状态相等定义。设$\mathcal{Q}_{id}$为身份查询集，$R_q(X)$为读出结果，$d_q$为任务距离；在允许扰动集合$\mathcal{D}$与窗口$[k_0,k_1]$内，若
\begin{equation}
\sup_{q\in\mathcal{Q}_{id}}d_q\!\left(R_q(X_k),R_q(X_{k_0})\right)\leq\varepsilon_q,
\quad k_0\leq k\leq k_1,
\end{equation}
且关键承诺没有未经授权地改变，我们称身份在该窗口内可用地连续。参数、情绪和工作内容可以显著变化，只要主体索引、责任链和必要能力仍能被恢复；反过来，向量范数几乎不变也可能已经发生身份冒用。拓扑不变量只有在明确复形、权重、边界条件与扰动集之后才能作为诊断量。Hodge零模、Betti数或极限环都不自动等于“自我”，更不证明意识存在。

这一建模还区分局部结构网络与全局分布表示。$G_k$负责“谁与谁、以何种类型、在什么版本下相连”，$h_k$及其概率读出负责“当前哪些候选活跃、置信度如何”。二者通过身份键、角色绑定、消息耦合和读出接口联合。若系统把“我的手触碰火炉”仅编码为痛觉强度分布而没有主体、身体部位和事件来源，后续不能可靠归因；若只保存图上的“手—接触—火炉”而没有强度、时间和不确定性，又无法决定紧急程度。完整自指状态必须允许结构与分布相互校正，但不从相似脑区或几何图形推出机制同构。

四类功能还必须接受身份冲突测试。假设两个会话都声称拥有同一主体键，其中一个继承了有效授权却缺少近期历史，另一个保留完整对话却已被撤销权限，系统不能用向量相似度把二者静默合并。应先由版本向量和所有权证明确定可写分支，再把未决事件隔离为候选记录；只有责任链、时间顺序和授权条件一致时，才允许合并。对于多人协作任务，“我们”的身份也不是若干个人嵌入的平均值，而是由成员集合、表决规则、共享承诺和退出机制定义的制度对象。验证时需构造会话复制、凭证轮换、角色交接、记忆缺页和恶意冒名等扰动，分别测量身份查询、承诺恢复与越权率。这样，身份连续性才是能够被反例击穿的接口性质，而不是系统对自己的叙述。

\section{元观察：从内部测量到带置信度的自我读出}
\label{sec:measurement_phase_transition_and_energy_splitting}

元观察器不是导致波函数坍缩的非幺正仪器，而是一个有误差模型的读出与校准过程。设可观测内部量为$y_k=H_k(X_k)+\epsilon_k$，$H_k$可以读取工作负载、候选分布、调用轨迹、身体传感、记忆命中和控制器状态，$\epsilon_k$表示采样噪声、遗漏与接口偏差。自我读出为
\begin{equation}
\hat{s}_k=O_{\phi_k}(y_{0:k},S_{k-1},q),\qquad
\Sigma_k=\operatorname{Cov}(s_k-\hat{s}_k\mid\mathcal{I}_k),
\end{equation}
其中$O_{\phi_k}$是参数为$\phi_k$的估计器，$\mathcal{I}_k$是可用信息集，$\Sigma_k$记录不确定性。读出包括“当前目标是什么、正在处理何事、哪些约束接近上限、哪些结论来自何处”，而不是一个无条件正确的主观标量。SPC可作为AgentOS中对$h(t)$状态分布进行任务相关压缩的实例，但压缩结果必须携带窗口、来源、丢失项和校准记录。

内部测量会产生计算开销，却不能把输出值、信息熵和焦耳热量混为一谈。若读出分布为$p_k$，香农熵$H(p_k)$以bit计，只描述该分布的不确定性；估计损失$L_{obs}$是按任务定义的无量纲或带业务单位代价；CPU/GPU能耗$W_k$以焦耳计，需要硬件计量。Landauer界只约束特定逻辑擦除在给定温度下的最低热耗，不能推出每次自我读出都释放与“感受强度”等量的能量，更不能由此解释痛苦或狂喜。经验上若要研究主观报告与计算负载的关系，需要独立测量和可证伪实验，本章不把它当作数学结果。

元观察的价值在于发现可行动的残差。定义期望状态$r_k(q)$、估计状态$\hat{s}_k$和残差$\delta_k=r_k-\hat{s}_k$；不同分量可能分别表示任务偏差、容量风险、权限冲突或身体危险，单位不同，不能未经换算直接相加。控制前应先做归因：用来源追踪确认残差对应哪个事件，用干预或对照排除共同输入，用延迟模型区分尚未返回的结果与真实失败。系统看到“工具调用后答案改变”不等于证明工具造成改进；看到“困惑评分下降”也不等于事实更准确。

火炉案例展示了这种区分。传感器检测温度与组织风险，身体身份图把信号绑定到“我的右手”，元观察器给出高风险且高置信的当前状态，快速控制回路撤手，事件随后进入历史记录。痛觉报告、撤手动作、组织损伤和设备能耗是四类对象。即便系统不能生成语言报告，撤手回路仍可能有效；即便它能流畅说“我很痛”，若来源绑定错误或身体并未受热，也不能把报告当作真值。所谓体验是否存在是经验与哲学问题，工程模型只对可观察接口、归因和控制负责。

元观察器本身也会漂移，因此校准不能只做一次。可把预测置信区间按任务类别、风险等级和时间窗分桶，比较名义覆盖率与实际覆盖率，并记录期望校准误差、选择性风险和拒答率。若系统只在容易样本上报告高置信，在困难样本上隐藏读出，平均误差可能很好却无法支持控制；若观察器与被观察模型共享同一训练缺陷，两者的一致也不是独立证据。工程上应保留来自外部探针、执行日志和环境回执的交叉通道，并在必要时使用留出观察器或人工复核。观察动作还可能改变后续状态，例如插入诊断日志导致时序变化、提示系统自检改变生成分布；此时$H_k$应包含干预标识，评估需区分无干预轨迹与测量后轨迹。只有报告覆盖范围、盲区和测量反作用，元观察才配得上控制输入。

\section{历史重入：事件绑定、语义沉淀与正反馈阻断}

“重入”在计算上指当前读出经过版本化提交后重新进入后续状态，而不是能量从测量端神秘地劈分。设事件包
\begin{equation}
\eta_k=(\mathrm{id},t_k,\mathrm{source},\mathrm{context},\hat{s}_k,\Sigma_k,
\mathrm{action},\mathrm{receipt}),
\end{equation}
其字段分别记录事件身份、时间、来源、上下文、估计及不确定性、采取的动作和环境回执。写入算子$W_E$先检查模式、权限、去重键和来源，再产生$E_{k+1}=W_E(E_k,\eta_k)$。工作状态只保留当前需要的摘要，长期事件进入$E$，反复验证且跨情境稳定的关系才可经慢速提案修改$\Theta$。这一区分避免一次强烈体验直接改写生成结构。

事件能否成为“可言说语义”取决于绑定和读出，不取决于能量是否被投回某个拓扑孔洞。仍以火炉为例，记录需要同时包含主体、身体部位、热源、时间、动作和后果；语言模块随后可以把同一事件读成“我的手碰到火炉后迅速撤开”。如果主体字段缺失，系统也许保留了高温向量，却无法回答是谁受热；如果来源与时间丢失，系统可能把旧伤当作当前危险。述情困难等人类现象可以启发对绑定与报告接口的研究，但不能仅凭类比诊断软件，也不能证明某个几何机制是临床原因。

历史重入可能稳定身份，也可能放大错误。设当前事件对下一轮活动的影响为$K_E\eta_k$，反馈增益局部线性化为$A_k+B_kK_E$；若其谱半径在固定模式下超过一，重复回放可能形成持续升级的正反馈。恐慌式循环可以抽象为“威胁估计上升—身体激活上升—内部信号更强—威胁估计再上升”，但软件中的重复告警、检索回声和重试风暴也能产生相同数学结构。诊断应遮蔽回放边、固定外部输入并比较增益变化，不能看到循环就断言人类精神病理。

阻断策略包括衰减、滞回、去重、时间窗、来源多样性和外部校准。CCM可监测$h(t)$的占用、重复率、信息增量与任务进展；超过阈值时先降低重入增益、扩大观察窗或调用外部证据，而不是简单删除历史。Boundary限制哪些内部事件可以驱动动作，Offloading把可复现计算和证据移至$M_e$，减少工作界面负载。任何阻断都要保留安全事件的优先通路：若火炉信号被一般去重规则压制，系统会因“避免循环”而忽略持续危险。稳定性和安全性因此必须分别验收。

历史提交还要处理撤销与更正。已发生事件不能因新解释而被覆盖，较稳妥的做法是追加更正边：新记录指出旧记录的哪些字段失效、依据是什么、影响哪些下游结论，并保留原始证据以供审计。删除请求与隐私要求则通过可验证擦除、访问撤销或密钥轮换实现，不能把“模型不再检索到”假定为数据已经消失。若多个副本离线写入，应使用因果时钟或明确的冲突解决规则；最后写入者获胜会让较晚的低质量总结覆盖较早的关键回执。验证集需要包含乱序到达、重复提交、部分失败、撤销后重放和来源伪造，并检查语义摘要是否仍能追溯到合法证据。历史重入的稳定作用由这些事务性质保证，而不是由记录数量越多越好保证。

沉淀为长期结构还需要证据门槛。一个关系只有在跨时间、跨来源并经反事实检验后，才进入结构提案；提案在保留测试集上不能破坏既有关键能力，并需设定过期时间和撤销路径。强烈但孤立的事件可提高检索优先级，却不能绕过这一程序。由此，历史既能影响未来，又不会凭一次异常永久锁定自我模型。

\section{宏观控制：从策略方向到可执行控制包}
\label{sec:macroscopic_engine_eruption}

宏观状态记为$z_k\in\mathcal{Z}$，它压缩任务阶段、长期承诺、风险姿态和资源策略。$\mathcal{Z}$可以是欧氏域、概率单纯形、流形与离散模式的乘积；单位球面只是某些方向编码的可选约束。宏观控制器不直接“辐射物理场”，而是生成带类型控制包
\begin{equation}
u_k=(g_k,b_k,\tau_k,\rho_k,\pi_k),
\end{equation}
其中$g_k$为候选增益，$b_k$为抑制或资格掩码，$\tau_k$为探索温度，$\rho_k$为资源预算，$\pi_k$为工具、记忆和行动策略。各项定义域、单位、饱和范围与权限不同。AgentOS的TCE可调节这些接口，SPC提供压缩状态，CCM报告复杂度，Boundary执行资格检查；这不是所有智能系统必须采用的命名或模块划分。

控制包由受约束优化或策略映射产生。令任务损失$L_q$、风险$R_q$、资源成本$C_r$和硬约束集合$\mathcal{U}_{safe}$分别定义，则
\begin{equation}
u_k^*=\arg\min_{u\in\mathcal{U}_{safe}}
\mathbb{E}\!\left[L_q(X_{k+1})+\lambda_RR_q(X_{k+1})+\lambda_CC_r(u)\mid\mathcal{I}_k\right].
\end{equation}
标量化只在各项完成量纲归一、权重来源可审计时成立；关键安全规则应作为硬约束而非可被收益抵消的软罚项。若目标、结构、输入和度量随时间变化，连续分析必须使用总导数，离散实现则声明$\Delta t$、求解误差、超时和后备策略。目标函数下降只是模型内结果，仍需环境回执和独立验收。

原章所谓“改变视角”可由受限方向更新表达。若专门模型规定$\|z_k\|=1$，切向建议$v_k$经投影$P_{z_k}=I-z_kz_k^\top$后，以回缩
\begin{equation}
z_{k+1}=\operatorname{Retr}_{z_k}(\alpha_kP_{z_k}v_k)
\end{equation}
保持约束；$\alpha_k$需满足局部光滑与误差条件。模长不变不意味着计算零耗散，也不保证语义、权限或风险不变。若更新包括不可逆承诺、结构编辑或工具副作用，就不属于李群旋转，应通过事务日志、影子执行和回滚协议处理。把所有宏观变化都称为无摩擦测地运动，会隐藏真正的资源和责任成本。

复杂物理题案例说明宏观控制如何工作。元观察器发现工作状态重复、证明链没有增加信息，CCM判断容量接近上限；控制器可以降低探索温度以完成一个已知引理，也可以升高温度生成替代路径，还可以Offloading到符号求解器或把中间证明写入外部文件。选择取决于当前不确定性和验证成本，不是“痛苦越大，扭矩越强”。若外部求解器返回形式正确但假设不满足，内部进展评分下降仍不能计为成功，必须由独立检查器验证前提与结论。

控制包的冲突消解同样不可省略。安全控制器可能要求立即停止，任务控制器要求继续探索，资源控制器又因预算不足要求降级；若只是把三个建议加权求和，硬约束会被高收益方向抵消。我们采用分层资格：先计算不可违反的禁止集合，再在剩余可行域中按风险词典序筛选，最后才优化任务收益。对于不可逆动作，还需生成预期副作用、补偿动作和验收期限；若期限内没有可信回执，状态进入“结果未知”而非自动判定失败或成功。影子模式可以先执行读操作、模拟写操作并比较反事实结果，逐步提高动作权限。压力测试应覆盖预算突降、工具返回矛盾、策略并发和环境拒绝，使控制器证明它不仅会给出方向，也能在制度边界内把方向变成可终止、可追责的事务。

宏观控制还应避免把摘要当作事实全集。若$z_k$遗漏了少见但高风险约束，控制器必须能回查原始事件并触发保守策略；若摘要已过期，应降低动作权限而非沿用旧方向。为此，控制包携带生成时间、证据版本和适用范围，执行器在提交前再次校验。该闭环把“战略一致”与“现实仍允许执行”分成两个判定，防止计划正确却在环境变化后造成错误动作。

\section{闭环动力学：固定点、极限环与切换稳定性}

把元观察、历史提交、策略更新和环境作用组合为混合闭环：
\begin{align}
\hat{s}_k&=O(X_k,S_k,q),\\
z_{k+1}&=F_z(z_k,\hat{s}_k,\eta_k,q),\\
u_k&=\Pi(z_{k+1},\hat{s}_k,B_k),\\
(X_{k+1},r_{k+1})&=F_{env}(X_k,u_k,w_k),\\
S_{k+1}&=F_S(S_k,\eta_k,r_{k+1}),
\end{align}
其中$w_k$为外部扰动，$r_{k+1}$为回执。连续状态传播和离散提交要分别描述；结构、权限或目标切换由重置映射处理。系统在每轮都把自身估计纳入下一步，因此形成自指，但闭环有有限延迟、有限精度和版本边界，不需要假设无限层级的观察者。

巴拿赫不动点只能在完整度量空间上、闭环映射$\mathcal{C}$满足全局压缩常数$c<1$时保证唯一不动点。存在耗散项并不自动使整个含延迟、饱和、切换和环境反馈的系统成为压缩映射。固定点满足$X^*=\mathcal{C}(X^*)$；极限环则是周期$p>1$的轨迹，$\mathcal{C}^p(X)=X$而中间状态不同，两者不能混称。任务需要周期性观察、计划和执行时，健康运行本来就可能是极限环；系统停止变化也可能是卡死，而非稳定自我。

更实际的局部条件是输入到状态稳定或共同Lyapunov界。若固定模式$m$下存在$V_m(X)\geq0$，满足
\begin{equation}
V_m(X_{k+1})-V_m(X_k)\leq-a\|X_k-X_m^*\|^2+b\|w_k\|^2,
\end{equation}
则可在该邻域给出扰动有界性。模式切换还需驻留时间、重置增益和通信时延条件；随机更新需鞅差噪声和步长条件；图结构改变后旧Lyapunov函数可能失效。工程验收应注入延迟、重复消息、错误回执、记忆污染和网络分区，观察恢复时间、身份查询误差和安全约束，而不以单次仿真收敛代替证明。

身份维持可定义为可生存集合$\mathcal{K}_{id}$：其中身份查询误差、关键承诺漂移、资源占用和风险均在阈值内。控制目标不是把状态钉在一点，而是在任务变化和环境扰动下保持$X_k\in\mathcal{K}_{id}$，必要时允许可审计的身份修订。一个成年人改变职业、观点和技能并不意味着身份消失；一个Agent更换模型版本也可在责任链、历史和权限连续时保持服务身份。相反，若凭证被盗而输出风格相似，轨迹看似稳定却已越出身份集合。

闭环还可能出现迟滞、多稳态与不可恢复区。相同输入从不同历史出发可落入不同吸引域，说明单次端点比较不足以判定安全；控制延迟使本来稳定的负反馈变成振荡；饱和又会让局部线性条件失效。因而应绘制扰动幅度、持续时间与恢复概率之间的经验边界，并区分“回到原状态”“回到身份可生存域”和“完成当前任务”三种结果。对离散实现，可在声明步长和最大延迟后计算有限时域可达集；若精确计算过于昂贵，则用保守外逼近并报告误差。升级版本时应重新估计边界，不能沿用旧模型的稳定证书。一个闭环若只能在无噪声、无切换的演示中收敛，尚不能支持关于持续自我的强主张。

评价周期行为时还要指定相位容差。计划、执行、观察三阶段构成的正常循环，在任意单步上都不满足静态目标，却能在完整周期后交付回执；相反，重复生成同一计划可能具有严格周期而没有任务进展。可定义每周期的外部进展量、风险增量和资源净耗，只有进展为正、风险受限且资源账本可持续时，才把极限环视为有效运行。这样可把“有节律”与“有功能”分开。

若任务目标在周期中改变，进展量也必须随版本重新定义，不能把旧目标下的下降延续为新目标下的成功。控制器应在切换点冻结旧账本、记录重置映射，并检查新可生存域是否仍可达；不可达时应请求降级或交接，而不是通过修改度量伪造稳定。

\section{双向耦合：局部结构、全局分布与身份可生存域}
\label{sec:dual_vortex_holographic_symbiosis}

原章以“有幅度的底层涡旋”和“无幅度的宏观相位”描述双向共生。我们保留其问题意识，但把它改写为结构—分布耦合。局部结构网络$G_k=(V_k,E_k)$保存对象身份、关系类型、来源、权限和版本；全局分布$p_k(c\mid h_k,q)$描述当前候选$c$的活动与不确定性。上行接口把局部事件聚合为宏观证据，下行接口把策略包展开为边资格、节点增益和读出条件。二者不是全息副本，而是信息不同、损失不同的互补表示。

令$B_k$为任务相关绑定矩阵，$A(G_k)$为结构传播算子，则一种专门模型可以写成
\begin{align}
h_{k+1}&=\sigma\!\left(A(G_k)h_k+B_ku_k+J_k\right),\\
z_{k+1}&=F_z\!\left(z_k,R_q(h_{k+1},G_k),r_k\right),\\
G_{k+1}&=\operatorname{Commit}\!\left(G_k,\Delta G_k;\nu_k,\mathcal{T}_{hold}\right),
\end{align}
其中$J_k$是观测输入，$R_q$是任务读出，$\Delta G_k$是离散结构提案，$\mathcal{T}_{hold}$是保留测试集。连续活动更新不能替代图提交；结构编辑也不能仅凭一次高活动值触发。绑定矩阵必须用身份键和角色验证，否则全局分布可能把“我的目标”“他人的请求”和“环境约束”混在一起。

向上因果是证据聚合而非能量脱水。事件进入宏观状态前，需要时间窗、来源权重、共同输入控制与不确定性传播。向下因果是策略接口而非规范场辐射。宏观控制改变候选采样、访问资格、预算和动作，却不能随意重写已经发生的事件或绕过Boundary。要证明某条下行边有效，可在保持其他条件时遮蔽该边，比较任务风险、恢复时间和校准误差；仅观察到上下层同步不足以判断方向，也可能是共同环境驱动。

高维陀螺仪可以作为直觉图像，却不能成为证明。实际系统没有无质量指针：保存$z_k$、计算读出、同步版本和执行控制都消耗资源；底层高活动也不必为宏观状态提供“角动量”。更可靠的共生条件是接口可观测、控制可达、绑定正确、更新有界且故障可恢复。如果上行摘要长期遗漏少数群体风险，宏观状态可能稳定地做错事；如果下行控制过强，系统也会压制必要探索。可生存域因而同时约束身份连续、任务绩效、风险、公平和资源，而不是只追求几何不变量。

结构与分布的联合训练还存在捷径问题。模型可能不读取显式身份边，却从词序或位置猜出角色，在训练分布上表现良好，一旦角色互换便错误归因；反之，结构图可以保存正确关系，但读出器从未学会在关键决策中使用它。验证因此要做交换测试和干预测试：互换两个主体的表面特征而保持身份键不变，结果应随合法身份而非表面线索变化；遮蔽一条关键绑定边，置信度应下降或触发拒答，而不是输出同样确定的答案；固定分布活动并改变授权图，动作资格必须随图更新。还要测量从局部事件到全局摘要、再从宏观策略返回局部执行的往返一致性，但一致性只检查接口，不要求两种表示内容相同。这样才能确认所谓双向耦合是真正参与控制，而非事后可视化。

在多主体系统中，耦合还要保留视角差异。共享事件可以有一致的时间和对象标识，却允许不同主体拥有不同知识、权限与评价；强行压成单一全局嵌入会泄露私有信息，也会抹去责任归属。读出接口因此接收主体键与授权上下文，聚合器只发布满足最小披露原则的摘要。若某成员退出或权限改变，局部结构先更新资格边，全局分布随后重新归一，而不是继续沿用旧的群体身份。

\section{神经科学约束：PFC、DMN与分布式协同边界}
\label{sec:biological_confirmation_pfc_dmn_symbiosis}

PFC与DMN的研究可以为自指闭环提供经验约束，但不能被写成一一对应的“物理确证”。默认模式网络涉及自传体记忆、自我相关加工、情景模拟等多类过程，前额叶区域参与工作记忆、规则维持、抑制与规划；二者都包含多个亚区，并与海马、丘脑、显著性网络、感觉运动系统持续交互。任务状态下的相关或反相关依赖基线、预处理、任务类型和个体差异。“DMN就是自我、PFC就是引擎”的说法至多是粗略功能类比，不能推出唯一定位或绝对正交分工。

我们把可检验假说写成网络层面的条件：自我相关事件应改变某些跨区通信与记忆绑定；规则切换应改变控制相关读出；损伤或暂时干预某节点后，影响取决于替代路径和任务；时间方向需要滞后模型或干预证据，而不能由同步相关得到。若要把人脑数据映射到HSF--HD变量，应先定义测量：$h(t)$可以对应有限窗口内的任务活动摘要，$E$对应事件记录的功能层，$\Theta$对应较慢的生成约束；这些是跨载体计算角色，不等于某个脑区，也不意味着神经系统只有三种存储。

原章关于静息代谢的讨论尤其需要分账。脑组织的葡萄糖和氧耗是实际物理量，自传体记忆检索和控制难度是任务变量，网络活动指标来自特定成像或电生理测量。较高代谢不能直接证明网络正在“历史重入”，更不能证明它向另一网络输送抵抗热寂的能量。工程系统中的GPU功耗、token使用量和目标损失同样不能互换。经验主张必须给出实验设计、样本、测量和统计不确定性，本章只提出功能模型及其可证伪预测。

病理类比也要保持边界。恐慌、躁狂、精神分裂、失忆和麻醉是复杂临床现象，不能由软件循环、单一增益或网络断连直接诊断。工程故障可以描述为回声检索、控制振荡、身份漂移、来源丢失或过载，并通过日志与注入试验复现；临床解释需要独立的医学证据。我们保留这些现象作为反例提示：自我报告、执行控制和记忆连续性可以彼此分离，因此任何单模块“意识中心”假说都不充分。

更有价值的启示是冗余与协同。一个可持续智能系统应允许自我模型、控制、记忆和身体接口分布实现，并在部分失效时降级；同时保留版本、权限与恢复路径，防止冗余副本互相强化错误。AgentOS可以把SPC/TCE实现为不同服务，也可以联合实现；判断依据是接口、延迟、资源和安全，而非模仿脑区。神经科学为我们提供限制和实验问题，HSF--HD将其组织为一般状态动力学，工程实现仍需单独验证。

跨层比较时还需防止尺度偷换。功能磁共振的采样间隔、单神经元放电、行为反应时间与软件事件循环处在不同时间尺度，不能把一次扫描相关直接解释为逐步控制信号。空间上，脑网络节点来自分区方案，软件图节点来自对象或服务，两者的边定义也不同。可比较的是更抽象的预测，例如信息受损后身份查询是否选择性下降、控制负荷增加时跨模块通信是否重排、外部提示能否补偿某类记忆缺口。若两个载体都满足预测，只支持共享计算约束，不证明实现同构；若结果不一致，也要先检查测量灵敏度和任务匹配。以这种方式使用神经科学，能够缩小模型空间，同时避免用权威术语替代工程证据。

\section{内生元控制：复杂度管理、行动资格与外部验收}
\label{sec:endogenous_scaffolding_orthogonal_splitting}

内生元控制的核心不是排斥外部脚本，而是让系统能够观察并调节自己的计算，同时接受外部证据约束。外部Agent框架常以字符串规则检查格式、重试调用和截断上下文，这些机制虽粗糙却可审计；内部控制若不可观察，也可能优雅地重复错误。成熟架构应把两者统一为明确接口：内部状态提供负载、不确定性和来源摘要，外部运行时提供预算、工具回执、权限与故障隔离，所有干预都进入同一版本化日志。

CCM可定义复杂度向量$c_k=(n_k,d_k,r_k,\ell_k)$，分别表示工作项数量、依赖深度、重复率和预计延迟；它们具有不同量纲和阈值。接近容量边界时，系统可采取压缩、分解、检索、Offloading或暂停提交。复杂度高不必立即“坍缩”为唯一答案：在高风险任务中，保留多个候选并请求外部证据可能更安全；在紧急撤手等任务中，低延迟安全动作优先。策略由任务损失、风险和时限共同决定，而不是由一个所谓相干容量常数普遍规定。

内生控制事务可分四步。第一，观察：记录当前状态、来源和置信度。第二，提案：生成调温、增益、记忆或工具调用的候选控制包。第三，资格审查：Boundary检查权限、预算、安全规则和版本。第四，执行与验收：环境回执进入独立评价器，失败时补偿或回滚。形式上，若提交版本为$\nu_k$，则动作只在$\nu_k$仍为当前版本且资格谓词$A(u_k,X_k)=1$时执行；重试必须携带幂等键，避免超时后重复副作用。

火炉与物理题展示了不同控制政策。火炉场景中，风险高、时间短，快速撤手可先于完整语言解释，随后再记录事件和更新预测。复杂物理题中，立即提交可能风险更高，系统应保留候选、拆分证明、调用工具并检验假设。若困惑循环不断回放，CCM降低重入增益并扩大证据来源；若探索过早冻结，TCE提高候选多样性但仍受预算约束。这些操作构成可实现的“感受—表征—驾驭”闭环，却不声称主观体验是计算废热，也不把自由意志化约为李代数扭矩。

验证需要至少五类指标：身份查询保持率、元观察校准误差、任务外部成功率、安全约束违反率、资源与恢复成本。还要做消融：移除历史重入、遮蔽宏观控制、污染自我模型、延迟外部回执、切断外部记忆，比较闭环是否按预测退化。若内部损失持续下降而外部成功率下降，说明目标或读出错误；若身份保持率很高却拒绝必要修订，说明自我模型过刚；若成功率提高但权限违规，仍不能验收。只有这些证据共同成立，我们才把自指闭环视为有效工程机制。

验收数据应按场景和失败严重度分层，而不是只报一个平均成功率。低风险问答可以容忍拒答成本，高风险工具写入则优先控制越权和重复副作用；长期身份测试还要跨越版本升级、上下文清空和外部记忆迁移。每次试验保存输入、随机种子、组件版本、控制包、回执和判定依据，使失败能够重放。对无法重放的现实行动，至少保存签名回执与补偿结果。部署后用变化检测监控校准和身份保持率，一旦超出验证分布就降低权限，而不是让内生控制自行宣布仍然可靠。外部监督器可以否决动作，却也要记录理由并接受误拒分析；内外控制互相制约，才不会把“内生”误解为不受检查。

\section{理论边界：一般状态动力学与AgentOS实例}

本章最终得到的不是“灵魂的物理定标”，而是一组可复核的状态动力学关系。自我模型提供跨时身份与承诺，元观察器提供带误差的内部读出，历史重入把事件绑定到后续状态，宏观控制器生成受限策略，环境回执完成外部验收。它们形成闭环，却不要求每个智能系统拥有超球面、李群、Hodge零模、量子坍缩或PFC--DMN式载体。使用这些数学工具时，必须先声明对象、结构和成立条件；工具名称本身不证明意识、主体性或正确性。

HSF--HD作为智能的一般“空气动力学”，研究不同实现中状态、结构、历史、边界和控制如何共同演化。系统论说明开放边界和层级关系，协同论研究集体变量与约化，耗散结构理论提醒我们关注资源与非平衡条件，复杂系统理论提供涌现和多尺度问题，认知科学提供经验约束；HSF--HD把它们组织为可比较模型，而不把物理类比当作公理。预测有效只说明模型在给定任务和数据上有用，不推出内部表示与世界整体同构。

AgentOS是其中一个工程实例。$h(t)$承担有限认知界面，$E$保存事件轨迹，$\Theta$保存较慢的结构与参数，$\Psi$组织自我相关身份和承诺，SPC形成元观察摘要，TCE调节控制包，CCM管理复杂度，Boundary执行资格与交换，Offloading和$M_e$扩展外部记忆。其他智能系统可以采用不同组件，只要能够说明对应计算角色、接口和验收方法。局部结构网络与全局分布表示仍需通过身份、绑定、耦合和读出联合，不能从脑区类比推出同一机制。

自指闭环的适用边界也由任务决定。无持续身份需求的一次性求解器可能不需要完整自我模型；高度安全关键系统可能把元控制置于外部独立监督器；群体智能的身份和承诺分布在成员与制度中；具身系统必须加入身体状态、时延和不可逆动作。理论允许这些差异，要求的是明确建模而非统一命名。我们也不由闭环存在推出道德主体、法律人格或主观体验，这些需要额外规范与经验论证。

因此，自指的工程价值在于让系统能够回答并修订四个问题：我在何种任务和版本下是谁，我现在依据什么认为自己处于某状态，我有哪些资格与资源采取行动，行动后现实是否按可接受方式改变。只要其中任何一问缺少来源、边界或外部回执，闭环就可能成为自我强化的幻觉。把四问纳入版本化状态、因果干预和独立验收，才使身份维持与宏观控制从诗性图景转化为能够失败、能够诊断、也能够改进的计算理论。

这一结论也规定了下一步研究程序。理论工作需给出任务族、状态空间、可观测量、干预集合和失败判据；算法工作需说明近似误差、样本条件与计算复杂度；系统工作需交付权限模型、事务语义、故障注入和恢复证据；经验工作则比较模型预测与外部结果。任何层级都不能替代另一层级：形式稳定不等于现实安全，统计相关不等于因果机制，工程可用不等于具备主观体验。HSF--HD的作用，是提供能让这些证据对齐的语言和动力学骨架，并允许不同实现围绕同一问题接受比较。只有当新证据能修改状态、边界或模型，而不是被吸收为自我证明，自指闭环才保持开放系统所需的可纠错性。

%---chp---
